\chapter{1977年上海卷（理）}

\begin{problem}
化简：\(\left(\frac{a}{a+b}-\frac{a^2}{a^2+2ab+b^2}\right)\div\left(\frac{a}{a+b}-\frac{a^2}{a^2-b^2}\right)\).

\end{problem}

\begin{problem}
计算：\(\frac{1}{2}\lg 25+\lg 2-\lg \sqrt{0.1}-\log_2 9\times\log_3 2\).

\end{problem}

\begin{problem}
\(\sqrt{-1}=\i\)，验算 \(\i\) 是否方程 \(2x^4+3x^3-3x^2+3x-5=0\) 的解.

\end{problem}

\begin{problem}
求证：\(\frac{\sin\left(\frac{\pi}{4}+\theta\right)}{\sin\left(\frac{\pi}{4}-\theta\right)}+\frac{\cos\left(\frac{\pi}{4}+\theta\right)}{\cos\left(\frac{\pi}{4}-\theta\right)}=\frac{2}{\cos 2\theta}\).

\end{problem}

\begin{problem}
在 \(\triangle ABC\) 中，\(\angle C\) 的平分线交 \(AB\) 于 \(D\)，过 \(D\) 作 \(BC\) 的平行线交 \(AC\) 于 \(E\)，已知 \(BC=a\)，\(AC=b\)，求 \(DE\) 的长.

\bitmapfigure[max width=6.0cm]{img/1977/shanghai_science/q2_fig1.png}

\end{problem}

\begin{problem}
已知圆 \(A\) 的直径为 \(2\sqrt{3}\)，圆 \(B\) 的直径为 \(4-2\sqrt{3}\)，圆 \(C\) 的直径为 \(2\)，圆 \(A\) 与圆 \(B\) 外切，圆 \(A\) 又与圆 \(C\) 外切，\(\angle A=60^\circ\)，求 \(BC\) 及 \(\angle C\).

\end{problem}

\begin{problem}
正六棱锥 \(V-ABCDEF\) 的高为 \(2\ \mathrm{cm}\)，底面边长为 \(2\ \mathrm{cm}\).

\begin{enumerate}
\item 按 \(1:1\) 画出它的二视图；
\item 求其侧面积；
\item 求它的侧棱和底面的夹角.
\end{enumerate}

\end{problem}

\begin{problem}
解不等式：\(\begin{cases}16-x^2\ge0,\\x^2-x-6>0\end{cases}\)，并在数轴上把它的解表示出来.

\end{problem}

\begin{problem}
已知两定点 \(A(-4,0)\)，\(B(4,0)\)，一动点 \(P(x,y)\) 与两定点 \(A\)，\(B\) 的连线 \(PA\)，\(PB\) 的斜率的乘积为 \(-\frac{1}{4}\). 求点 \(P\) 的轨迹方程，并把它化为标准方程，指出是什么曲线.

\end{problem}

\begin{problem}
等腰梯形的周长为 \(60\)，底角为 \(60^\circ\)，问这梯形各边长为多少时，面积最大？

\bitmapfigure[max width=6.0cm]{img/1977/shanghai_science/q7_fig1.png}

\end{problem}

\begin{problem}
当 \(k\) 为何值时，方程组 \(\begin{cases}x-\sqrt{y-2}=0,\\kx-y-2k-10=0\end{cases}\) 有两组相同的解？并求出它的解.

\end{problem}


\section{附加题}

\begin{problem}
如图所示，半圆 \(O\) 的直径为 \(2\)，\(A\) 为半圆直径的延长线上的一点，且 \(OA=2\)，\(B\) 为半圆上任一点，以 \(AB\) 为边作等边 \(\triangle ABC\)，问 \(B\) 在什么地方时，四边形 \(OACB\) 的面积最大？并求出这个面积的最大值.

\bitmapfigure[max width=6.0cm]{img/1977/shanghai_science/q9_fig1.png}

\end{problem}

\begin{problem}
已知曲线 \(y=x^2-2x+3\) 与直线 \(y=x+3\) 相交于点 \(P(0,3)\)，\(Q(3,6)\) 两点.

\begin{enumerate}
\item 分别求出曲线在交点的切线的斜率；
\item 求出曲线与直线所围成的图形的面积.
\end{enumerate}
\end{problem}
